% GENERATED FILE. Edit canonical lessons, then run npm run generate:books.
% canonical-source-hash: fnv1a64:64bf4dd9b068b936
% canonical-lessons: RU-C356-nesmotrya-na, RU-C356-tem-ne-menee, RU-C356-vprochem, RU-C356-s-odnoy-storony, RU-C356-s-drugoy-storony

\chapter{Despite everything}
\label{ch:ru-despite-everything}
\hlchaptermodality{356}

\begin{chapteropening}
\textbf{By the end of this chapter:} \emph{I can say despite, nevertheless and however, and set one side of a question against the other.}

Complete the last lesson of chapter 356: I can say despite, nevertheless and however, and set one side of a question against the other.
\end{chapteropening}

\section[\texorpdfstring{\ru{несмотря} \ru{на} (nesmotryá na)}{несмотря на (nesmotryá na)}]{\ru{несмотря} \ru{на} (nesmotryá na) — despite, in spite of}

\label{lesson:RU-C356-nesmotrya-na}



\begin{quote}
Before the new one: say the Russian for inevitable, then the Russian for to arise.
\end{quote}

\subsection*{You'll want to know: \ru{несмотря} \ru{на}}
\textbf{\ru{несмотря} \ru{на}} (\emph{nesmotryá na}) — \textquotedblleft{}despite, in spite of\textquotedblright{}.

\begin{quote}
\ru{Несмотря} \ru{на} \ru{дождь}, \ru{мы} \ru{пошли} \ru{в} \ru{парк}.
\end{quote}

\begin{quote}
\emph{Nesmotryá na dozhd, my poshlí v park.}
\end{quote}

\begin{quote}
Despite the rain, we went to the park.
\end{quote}

It takes the accusative: \emph{\ru{несмотря} \ru{на} \ru{погоду}}. Before a clause, \emph{\ru{несмотря} \ru{на} \ru{то}, \ru{что}…} — although. \emph{\ru{Несмотря}} is written as one word.

\begin{cousinweb}
\textbf{\ru{несмотря} \ru{на}} is built from \emph{\ru{не}} \textquotedblleft{}not\textquotedblright{} and \emph{\ru{смотря}}, the gerund of \emph{\ru{смотреть}} \textquotedblleft{}to look\textquotedblright{}, from Proto-Slavic \emph{*motriti} \textquotedblleft{}to watch\textquotedblright{}: not looking at the obstacle.
\end{cousinweb}

\begin{grammarlens}[title={What you've built}]
One more word for telling, explaining and describing at B1.
\end{grammarlens}

\subsection*{Your turn}
\begin{itemize}
\raggedright
  \item \emph{Say it:} \emph{nesmotryá na}
  \item \emph{Say it:} \emph{nesmotryá na}, once more
  \item \emph{Say it:} a sentence of your own with \emph{nesmotryá na}
  \item \emph{Recall:} say the Russian for inevitable, then the Russian for to arise, then say \emph{nesmotryá na} again
\end{itemize}

\begin{tcolorbox}[breakable,colback=teal!4,colframe=teal!35!black,title={Before you move on}]
What does \ru{несмотря} \ru{на} mean? (\textquotedblleft{}Despite, in spite of\textquotedblright{}.) Say it once more.
\end{tcolorbox}
\section[\texorpdfstring{\ru{тем} \ru{не} \ru{менее} (tem ne méneye)}{тем не менее (tem ne méneye)}]{\ru{тем} \ru{не} \ru{менее} (tem ne méneye) — nevertheless}

\label{lesson:RU-C356-tem-ne-menee}



\begin{quote}
Before the new one: say the Russian for to arise, then the Russian for despite.
\end{quote}

\subsection*{You'll want to know: \ru{тем} \ru{не} \ru{менее}}
\textbf{\ru{тем} \ru{не} \ru{менее}} (\emph{tem ne méneye}) — \textquotedblleft{}nevertheless\textquotedblright{}.

\begin{quote}
\ru{Было} \ru{трудно}, \ru{тем} \ru{не} \ru{менее} \ru{мы} \ru{всё} \ru{сделали}.
\end{quote}

\begin{quote}
\emph{Býlo trúdno, tem ne méneye my vsyo sdélali.}
\end{quote}

\begin{quote}
It was hard; nevertheless, we did everything.
\end{quote}

Formal. It may open the clause or follow \emph{\ru{но}}: \emph{\ru{но} \ru{тем} \ru{не} \ru{менее}}.

\begin{cousinweb}
\textbf{\ru{тем} \ru{не} \ru{менее}} is built from \emph{\ru{тем}} \textquotedblleft{}by that\textquotedblright{}, \emph{\ru{не}} \textquotedblleft{}not\textquotedblright{} and \emph{\ru{менее}} \textquotedblleft{}less\textquotedblright{}, the old comparative: not the less for that.
\end{cousinweb}

\begin{grammarlens}[title={What you've built}]
One more word for telling, explaining and describing at B1.
\end{grammarlens}

\subsection*{Your turn}
\begin{itemize}
\raggedright
  \item \emph{Say it:} \emph{tem ne méneye}
  \item \emph{Say it:} \emph{tem ne méneye}, once more
  \item \emph{Say it:} a sentence of your own with \emph{tem ne méneye}
  \item \emph{Recall:} say the Russian for to arise, then the Russian for despite, then say \emph{tem ne méneye} again
\end{itemize}

\begin{tcolorbox}[breakable,colback=teal!4,colframe=teal!35!black,title={Before you move on}]
What does \ru{тем} \ru{не} \ru{менее} mean? (\textquotedblleft{}Nevertheless\textquotedblright{}.) Say it once more.
\end{tcolorbox}
\section[\texorpdfstring{\ru{впрочем} (vpróchem)}{впрочем (vpróchem)}]{\ru{впрочем} (vpróchem) — anyway, mind you}

\label{lesson:RU-C356-vprochem}



\begin{quote}
Before the new one: say the Russian for despite, then the Russian for nevertheless.
\end{quote}

\subsection*{You'll want to know: \ru{впрочем}}
\textbf{\ru{впрочем}} (\emph{vpróchem}) — \textquotedblleft{}anyway, mind you\textquotedblright{}.

\begin{quote}
\ru{Я} \ru{не} \ru{люблю} \ru{мясо}, \ru{впрочем}, \ru{сегодня} \ru{попробую}.
\end{quote}

\begin{quote}
\emph{Ya ne lyublyú myáso, vpróchem, segódnya popróbuyu.}
\end{quote}

\begin{quote}
I do not like meat; still, today I will try some.
\end{quote}

It softens or partly takes back what was said: \emph{\ru{Впрочем}, \ru{это} \ru{неважно}} — anyway, it does not matter.

\begin{cousinweb}
\textbf{\ru{впрочем}} is fused from \emph{\ru{в}} \textquotedblleft{}in\textquotedblright{} and \emph{\ru{прочем}}, a form of \emph{\ru{прочий}} \textquotedblleft{}other, the rest\textquotedblright{}: as for the rest.
\end{cousinweb}

\begin{grammarlens}[title={What you've built}]
One more word for telling, explaining and describing at B1.
\end{grammarlens}

\subsection*{Your turn}
\begin{itemize}
\raggedright
  \item \emph{Say it:} \emph{vpróchem}
  \item \emph{Say it:} \emph{vpróchem}, once more
  \item \emph{Say it:} a sentence of your own with \emph{vpróchem}
  \item \emph{Recall:} say the Russian for despite, then the Russian for nevertheless, then say \emph{vpróchem} again
\end{itemize}

\begin{tcolorbox}[breakable,colback=teal!4,colframe=teal!35!black,title={Before you move on}]
What does \ru{впрочем} mean? (\textquotedblleft{}Anyway, mind you\textquotedblright{}.) Say it once more.
\end{tcolorbox}
\section[\texorpdfstring{\ru{с} \ru{одной} \ru{стороны} (s odnóy storoný)}{с одной стороны (s odnóy storoný)}]{\ru{с} \ru{одной} \ru{стороны} (s odnóy storoný) — on the one hand}

\label{lesson:RU-C356-s-odnoy-storony}



\begin{quote}
Before the new one: say the Russian for nevertheless, then the Russian for anyway.
\end{quote}

\subsection*{You'll want to know: \ru{с} \ru{одной} \ru{стороны}}
\textbf{\ru{с} \ru{одной} \ru{стороны}} (\emph{s odnóy storoný}) — \textquotedblleft{}on the one hand\textquotedblright{}.

\begin{quote}
\ru{С} \ru{одной} \ru{стороны}, \ru{работа} \ru{интересная}.
\end{quote}

\begin{quote}
\emph{S odnóy storoný, rabóta interésnaya.}
\end{quote}

\begin{quote}
On the one hand, the work is interesting.
\end{quote}

It sets up a first view, and a second view follows. \emph{\ru{Сторона}} is a side.

\begin{cousinweb}
\textbf{\ru{с} \ru{одной} \ru{стороны}} is \textquotedblleft{}from one side\textquotedblright{}: \emph{\ru{сторона}} \textquotedblleft{}a side\textquotedblright{} is inherited from Proto-Slavic \emph{*storna}.
\end{cousinweb}

\begin{grammarlens}[title={What you've built}]
One more word for telling, explaining and describing at B1.
\end{grammarlens}

\subsection*{Your turn}
\begin{itemize}
\raggedright
  \item \emph{Say it:} \emph{s odnóy storoný}
  \item \emph{Say it:} \emph{s odnóy storoný}, once more
  \item \emph{Say it:} a sentence of your own with \emph{s odnóy storoný}
  \item \emph{Recall:} say the Russian for nevertheless, then the Russian for anyway, then say \emph{s odnóy storoný} again
\end{itemize}

\begin{tcolorbox}[breakable,colback=teal!4,colframe=teal!35!black,title={Before you move on}]
What does \ru{с} \ru{одной} \ru{стороны} mean? (\textquotedblleft{}On the one hand\textquotedblright{}.) Say it once more.
\end{tcolorbox}
\section[\texorpdfstring{\ru{с} \ru{другой} \ru{стороны} (s drugóy storoný)}{с другой стороны (s drugóy storoný)}]{\ru{с} \ru{другой} \ru{стороны} (s drugóy storoný) — on the other hand}

\label{lesson:RU-C356-s-drugoy-storony}



\begin{quote}
Before the new one: say the Russian for anyway, then the Russian for on the one hand.
\end{quote}

\subsection*{You'll want to know: \ru{с} \ru{другой} \ru{стороны}}
\textbf{\ru{с} \ru{другой} \ru{стороны}} (\emph{s drugóy storoný}) — \textquotedblleft{}on the other hand\textquotedblright{}.

\begin{quote}
\ru{С} \ru{другой} \ru{стороны}, \ru{зарплата} \ru{маленькая}.
\end{quote}

\begin{quote}
\emph{S drugóy storoný, zarpláta málenkaya.}
\end{quote}

\begin{quote}
On the other hand, the pay is low.
\end{quote}

It follows \emph{\ru{с} \ru{одной} \ru{стороны}}, or stands alone to bring in a second view.

\begin{cousinweb}
\textbf{\ru{с} \ru{другой} \ru{стороны}} is \textquotedblleft{}from the other side\textquotedblright{}: \emph{\ru{другой}} \textquotedblleft{}other\textquotedblright{} and \emph{\ru{сторона}} \textquotedblleft{}a side\textquotedblright{}, inherited from Proto-Slavic \emph{*storna}.
\end{cousinweb}

\begin{grammarlens}[title={What you've built}]
One more word for telling, explaining and describing at B1.
\end{grammarlens}

\subsection*{Your turn}
\begin{itemize}
\raggedright
  \item \emph{Say it:} \emph{s drugóy storoný}
  \item \emph{Say it:} \emph{s drugóy storoný}, once more
  \item \emph{Say it:} a sentence of your own with \emph{s drugóy storoný}
  \item \emph{Recall:} say the Russian for anyway, then the Russian for on the one hand, then say \emph{s drugóy storoný} again
\end{itemize}

\begin{tcolorbox}[breakable,colback=teal!4,colframe=teal!35!black,title={Before you move on}]
What does \ru{с} \ru{другой} \ru{стороны} mean? (\textquotedblleft{}On the other hand\textquotedblright{}.) Say it once more.
\end{tcolorbox}
